\clearpage
\fancyhead[L]{\small COMET.}
\hypertarget{native-456-comet}{}
\begin{center}
{\large COMET.}
\end{center}

\begin{center}
\rule{6cm}{0.4pt}\\
Göttingische gelehrte Anzeigen. No. 57. pp. 561..565. 10 April 1826.\\
\rule{6cm}{0.4pt}
\end{center}

Our journal gave the first public account in 1819, in nos. 28 and 83 [pp. 417 and 420 of this volume], of the important discovery that made known to us a comet with an orbital period of only three years. It was named \SourceLeaf{457}\textsc{Encke}’s comet after the astronomer who contributed so much to its theory. The account ended with the wish that \emph{this most remarkable discovery might spur friends of astronomy to redouble their zeal in searching for comets, since everything suggested that it would be only the beginning of an immense harvest gradually coming to ripeness.} The 14 new comets observed and calculated since then show how far that \emph{wish} has been fulfilled. With a high degree of probability we can now also report the first confirmation of the \emph{conjecture}. On 28 February of this year, Captain \textsc{Biela}, in Josephsstadt in Bohemia, discovered a small comet in the constellation Aries, invisible to the naked eye. Following his announcement it was observed at several other places and has so far been observed twice at our observatory by Professor \textsc{Harding}. A few days later, Mr.\ \textsc{Gambard} in Marseille independently discovered the same comet. Given the large number of new comets found in recent years, we would give this one no special mention in our journal if the very first calculations of its orbit by several different people, based on different observations as well, had not yielded a most remarkable result. Court Councillor \textsc{Gauss} received the parabolic elements from five astronomers, Captain von \textsc{Biela} himself, Mr.\ \textsc{Clausen} in Altona, Professor \textsc{Encke} in Berlin, Dr.\ \textsc{Olbers} in Bremen and Professor \textsc{Schwerdt} in Speyer. We bring them together here:

\begin{center}
\begingroup\small\setlength{\tabcolsep}{5pt}
\begin{tabular}{rlll}
\multicolumn{4}{l}{Passage through perihelion} \\
1826. March & 15.45653 & v. B. & \\
 & 15.75422 & Cl. & \\
 & 18.2432 & Göttingen T. & E. \\
 & 18.33773 & --- & O. \\
 & 17.620 & Sch. & \\
\end{tabular}
\par\medskip
\begin{tabular}{ll}
\multicolumn{2}{l}{Longitude of perihelion} \\
$96^{\circ}\quad27'\quad33$ & v. B. \\
$95\quad48\quad15$ & Cl. \\
$107\quad55\quad31$ & E. \\
$112\quad39\quad24$ & O \\
$100\quad51\quad57$ & Sch. \\
\end{tabular}
\par\medskip
\begin{tabular}{ll}
\multicolumn{2}{l}{Ascending node} \\
$245^{\circ}\quad57'\quad0''$ & v. B. \\
$245\quad16\quad40$ & Cl. \\
$248\quad41\quad59$ & E. \\
$250\quad44\quad21$ & O. \\
$246\quad34\quad52$ & Sch. \\
\end{tabular}
\par\medskip
\begin{tabular}{ll}
\multicolumn{2}{l}{Inclination of the orbit} \\
$15^{\circ}\quad28'\quad23''$ & v. B. \\
$15\quad45\quad19$ & Cl. \\
$13\quad12\quad37$ & E. \\
$12\quad18\quad31$ & O. \\
$15\quad7\quad5$ & Sch. \\
\end{tabular}
\par\medskip
\begin{tabular}{ll}
\multicolumn{2}{l}{Logarithm of the least distance} \\
0.00506 & v. B. \\
0.0114369 & Cl. \\
9.96474 & E. \\
9.94460 & O. \\
9.99565 & Sch. \\
\end{tabular}
\endgroup
\end{center}
\noindent Prograde motion.

The differences between these results are small when one considers how scant the underlying data are, and it should soon be possible to reconcile them.

What immediately strikes us is how closely these elements resemble those of the second comet of 1805 in every respect. That resemblance alone makes it exceedingly probable that the two comets are the same. The circumstances of the comet of 1805 make this even more probable. Court Councillor \textsc{Gauss}, who had investigated its appearance at great length, found that the \SourceLeaf{458}observations pointed very strikingly to a short orbital period, although they did not suffice to determine it with any precision from this appearance alone. The observations were best reconciled with an orbital period of $4\frac{3}{4}$ years. According to those investigations, a period one year shorter or several years longer could also have been made to agree very well with the observations, and much better than parabolic motion. Moreover, the elements of the comet of 1805 had shown a close resemblance to those of the comet of 1772, leading to the conjecture that the two might be the same. The difference was still so considerable, however, that some astronomers regarded their identity as improbable or even regarded their being distinct as certain. Court Councillor \textsc{Gauss} showed even then that this judgment was insufficiently founded and that the possibility of fully explaining the difference was by no means excluded. Although he has not yet undertaken the lengthy calculations needed to decide the question, he has continued to think it probable that the two comets were the same. In these circumstances it is most remarkable that the time from the perihelion passage of the comet of 1772 to that of 1805 and the time from the passage of the comet of 1805 to that of 1826 are nearly in the ratio of the numbers 5 and 3.

We may therefore conjecture with high probability that all three comets are the same, that from 1772 to 1826 this comet made eight revolutions in all, each lasting an average of 6 years and 9 months, and that the differences in the elements of 1772 and 1805 that had made some astronomers doubt the conjectured identity were chiefly due to perturbations by Jupiter, to which the comet came very close in the intervening time. Without this last circumstance it would for the time being be just as probable that the comet had made 10 and 6 revolutions in the meantime, thus 16 in all, each of $3\frac{3}{8}$ years.

It is much to be hoped that the comet discovered by Mr.\ von \textsc{Biela}, observations of which were interrupted by moonlight in the middle of March, can be observed well again several times somewhere once the moonlight is gone. Even if this does not succeed, the observations already obtained will suffice, once properly treated, to decide whether these three comets are the same, as we are so irresistibly led to conjecture.